\documentclass[11pt,a4paper]{article}

\usepackage[utf8]{inputenc}
\usepackage[T1]{fontenc}
\usepackage[margin=2.5cm]{geometry}
\usepackage{hyperref}
\usepackage{xcolor}
\usepackage{listings}
\usepackage{enumitem}
\usepackage{booktabs}
\usepackage{longtable}
\usepackage{amsmath,amssymb}
\usepackage{fancyhdr}
\usepackage{titlesec}

\definecolor{isls-blue}{HTML}{1a5276}
\definecolor{isls-orange}{HTML}{e67e22}
\definecolor{isls-green}{HTML}{27ae60}
\definecolor{isls-red}{HTML}{c0392b}
\definecolor{isls-gray}{HTML}{7f8c8d}
\definecolor{code-bg}{HTML}{f7f9fb}

\lstdefinestyle{rust}{
  language=C,
  basicstyle=\ttfamily\small,
  keywordstyle=\color{isls-blue}\bfseries,
  commentstyle=\color{isls-gray}\itshape,
  stringstyle=\color{isls-green},
  backgroundcolor=\color{code-bg},
  frame=single,
  rulecolor=\color{isls-gray!30},
  breaklines=true,
  showstringspaces=false,
  tabsize=4,
  morekeywords={fn,let,mut,pub,struct,impl,enum,use,mod,self,Self,
    crate,super,where,trait,for,in,if,else,match,return,async,await,
    Ok,Err,Some,None,Vec,String,HashMap,BTreeMap,Result,Option,
    f64,f32,usize,bool,u32,u64,i64,Box,dyn,Arc},
}
\lstdefinestyle{bash}{
  language=bash,
  basicstyle=\ttfamily\small,
  backgroundcolor=\color{code-bg},
  frame=single,
  rulecolor=\color{isls-gray!30},
  breaklines=true,
  showstringspaces=false,
}

\pagestyle{fancy}
\fancyhf{}
\fancyhead[L]{\small\textcolor{isls-gray}{ISLS I1 Spec --- Infogenetik-Kern}}
\fancyhead[R]{\small\textcolor{isls-gray}{v1.0.0 --- \today}}
\fancyfoot[C]{\thepage}

\titleformat{\section}
  {\Large\bfseries\color{isls-blue}}{\thesection}{1em}{}
\titleformat{\subsection}
  {\large\bfseries\color{isls-blue!80}}{\thesubsection}{1em}{}
\titleformat{\subsubsection}
  {\normalsize\bfseries\color{isls-blue!60}}{\thesubsubsection}{1em}{}

\hypersetup{
  colorlinks=true,
  linkcolor=isls-blue,
  urlcolor=isls-orange,
}

\begin{document}

\begin{titlepage}
\centering
\vspace*{2.5cm}

{\Huge\bfseries\color{isls-blue} ISLS I1 Specification}\\[0.5cm]
{\LARGE\color{isls-orange} Infogenetik-Kern}\\[0.3cm]
{\large\color{isls-gray} 5D Codematrix, Multiskalen-Gate,\\
Resonite-Observables, Fitness-R\"uckkopplung}\\[2cm]

{\large
\begin{tabular}{rl}
\textbf{System:}       & ISLS --- Intelligent Semantic Ledger Substrate \\
\textbf{Version:}      & v6.0 (I1) \\
\textbf{Author:}       & Sebastian Klemm \\
\textbf{Date:}         & \today \\
\textbf{Status:}       & Specification \\
\textbf{Depends:}      & S1 (complete), M1 (complete) \\
\textbf{Branch:}       & \texttt{claude/implement-isls-i1-spec} \\
\textbf{References:}   & 5D Language v1.0.7, Infogenetik: Spektrale \\
                        & Zielsignaturen, TIMES Universal Layer, \\
                        & PSP Schnitt 1.2 SRI-MERKABA v3.0 \\
\end{tabular}
}

\vfill

{\small\color{isls-gray}
Mithras Substructure University\\
Repository: \texttt{https://github.com/LashSesh/graphity}\\
License: MIT
}
\end{titlepage}

\tableofcontents
\newpage


% ══════════════════════════════════════════════════════════════════════════════
\section{Executive Summary}
% ══════════════════════════════════════════════════════════════════════════════

After D8, M1, and S1, ISLS generates code, measures consensus 
(Ophanim Swarm), and feeds a self-learning norm pipeline. But 
the measurement is crude: binary (compiles/doesn't) at project 
level, Jaccard similarity at file level. The system has no 
\emph{dimensional} understanding of code quality.

I1 installs the \textbf{Infogenetik-Kern}: a formal measurement 
and feedback layer based on the 5D Codematrix (from the 5D 
Language Specification), Resonite observables (from the TIMES 
Universal Layer), Multiskalen-Gating (from the Infogenetik 
Spektrale Zielsignaturen), and evolutionary fitness tracking.

This is the nervous system upgrade. The pipeline stops being 
blind between ``LLM returned something'' and ``cargo check 
says yes/no.'' Every generated file gets a 5D fingerprint. 
Every layer gets a gate. Every generation feeds fitness back 
to the norms that caused it.

\subsection{What I1 Installs}

\begin{enumerate}[nosep]
\item \textbf{5D Codematrix:} Every generated Rust file is 
      measured on 5 axes: Relationality, Functional Cohesion, 
      Topology, Symmetry, Entropy. Result: a \texttt{[f64; 5]} 
      fingerprint per file.
\item \textbf{Resonites:} Formal observable primitives extracted 
      from code --- functions, structs, imports, layers, 
      connectivity --- each typed and measured. The atoms of 
      the infogenetik.
\item \textbf{Multiskalen-Gate:} Three gate levels inserted 
      into the StagedClosure pipeline: Mikro (per file), 
      Meso (per module), Makro (cargo check). Errors caught 
      earlier $=$ fewer Coagula cycles.
\item \textbf{Fitness-R\"uckkopplung:} Every norm that was 
      activated in a generation receives a fitness update 
      based on the outcome. Norms that consistently lead to 
      clean compiles get higher confidence. Norms that 
      consistently fail get suppressed.
\item \textbf{Kontraktivit\"ats-Monitor:} Each Coagula cycle 
      logs error count. If $\text{errors}(k+1) \geq 
      \text{errors}(k)$: non-contractive flag. The system 
      knows when it is not converging.
\end{enumerate}


% ══════════════════════════════════════════════════════════════════════════════
\section{The 5D Codematrix}
\label{sec:codematrix}
% ══════════════════════════════════════════════════════════════════════════════

\subsection{The Five Axes}

From the 5D Language Specification v1.0.7, adapted for 
code-generation context:

\begin{longtable}{clp{7cm}}
\toprule
\textbf{Axis} & \textbf{Symbol} & \textbf{What It Measures} \\
\midrule
Relationality & $R$ 
  & How connected is this file to others? Import count, 
    FK references, cross-module dependencies. High $R$: 
    hub file (main.rs, mod.rs). Low $R$: isolated utility. \\
Functional Cohesion & $F$ 
  & How focused is this file? Does it do one thing (high $F$) 
    or many unrelated things (low $F$)? Measured by the ratio 
    of internal references to external references. \\
Topology & $T$ 
  & Does the file's structure match its expected position in 
    the HDAG layer? A Layer-4 query file should look like a 
    query file, not like an API handler. \\
Symmetry & $S$ 
  & How consistent is naming? Do function names follow the 
    \texttt{get\_/list\_/create\_/update\_/delete\_} pattern? 
    Are types consistently used? \\
Entropy & $E$ 
  & Code complexity. Low entropy $=$ boilerplate 
    (models, config). High entropy $=$ algorithmic content 
    (queries, business logic). Neither extreme is bad --- but 
    it should match expectation for the file type. \\
\bottomrule
\end{longtable}

\subsection{Measurement Functions}

\begin{lstlisting}[style=rust]
/// 5D Codematrix point for a single generated file.
#[derive(Debug, Clone, Copy, Serialize, Deserialize)]
pub struct Codematrix {
    pub r: f64,  // Relationality [0, 1]
    pub f: f64,  // Functional Cohesion [0, 1]
    pub t: f64,  // Topology conformance [0, 1]
    pub s: f64,  // Symmetry [0, 1]
    pub e: f64,  // Entropy [0, 1]
}

impl Codematrix {
    /// Resonance product (from PSP Schnitt 1.2 Eq. 1).
    pub fn resonance(&self) -> f64 {
        (self.r * self.f * self.t * self.s * self.e).powf(0.2)
    }
    
    /// Distance between two codematrix points.
    pub fn distance(&self, other: &Codematrix) -> f64 {
        ((self.r - other.r).powi(2)
         + (self.f - other.f).powi(2)
         + (self.t - other.t).powi(2)
         + (self.s - other.s).powi(2)
         + (self.e - other.e).powi(2))
        .sqrt()
    }
}
\end{lstlisting}

\subsection{Computing Each Axis}

All computations use \texttt{isls\_reader::parse\_string()} 
to extract structure. No heuristics --- deterministic formulas.

\subsubsection{R: Relationality}

\begin{equation}
R = \min\!\left(1.0,\;\frac{|\text{imports}| + |\text{FKs}|}{R_\text{max}}\right)
\end{equation}

where $R_\text{max} = 15$ (calibration constant: a file with 
15+ imports/FKs is maximally relational).

\subsubsection{F: Functional Cohesion}

\begin{equation}
F = \frac{|\text{internal\_calls}|}{|\text{internal\_calls}| + |\text{external\_calls}| + \epsilon}
\end{equation}

Internal calls: function calls to functions defined in the 
same file. External calls: function calls to imported functions. 
A file where everything calls everything else internally: $F \to 1$.

Approximation (when call-graph is not available): 
$F = 1 - \frac{|\text{imports}|}{|\text{imports}| + |\text{functions}| + \epsilon}$.

\subsubsection{T: Topology Conformance}

\begin{equation}
T = \text{sim}(\text{resonites}(f), \text{expected\_resonites}(\text{layer}(f)))
\end{equation}

where \texttt{expected\_resonites(layer)} is the norm-derived 
expectation for files in that HDAG layer. Example: Layer 4 
(database queries) should have: \texttt{sqlx::query\_as}, 
\texttt{PgPool}, \texttt{AppError}, 5 CRUD functions. 
$T$ measures how well the actual file matches this expectation.

\subsubsection{S: Symmetry}

\begin{equation}
S = \frac{|\text{conforming\_names}|}{|\text{all\_names}|}
\end{equation}

Conforming names: functions that follow the ISLS naming 
convention (\texttt{get\_\{e\}}, \texttt{list\_\{es\}}, 
\texttt{create\_\{e\}}, \texttt{update\_\{e\}}, 
\texttt{delete\_\{e\}}) and structs that follow PascalCase 
with matching snake\_case modules.

\subsubsection{E: Entropy}

\begin{equation}
E = \frac{H(\text{token\_distribution})}{H_\text{max}}
\end{equation}

where $H$ is Shannon entropy over the distribution of 
Rust keywords and identifiers. Normalized by $H_\text{max} = 
\log_2(|\text{vocabulary}|)$.

Low $E$ ($< 0.3$): highly repetitive (struct definitions, config). 
High $E$ ($> 0.7$): complex logic (algorithms, queries). 
The measurement does not judge --- it characterizes.


% ══════════════════════════════════════════════════════════════════════════════
\section{Resonites: Observable Primitives}
\label{sec:resonites}
% ══════════════════════════════════════════════════════════════════════════════

From the TIMES Universal Layer Specification, operationalized 
for Rust code:

\begin{lstlisting}[style=rust]
/// Atomic observable extracted from Rust source code.
/// The building blocks of the Infogenetik.
#[derive(Debug, Clone, Serialize, Deserialize, Hash, Eq, PartialEq)]
pub enum Resonite {
    /// A function signature.
    Fn {
        name: String,
        arity: usize,
        return_type: String,
        is_pub: bool,
        is_async: bool,
    },
    /// A struct or enum definition.
    Type {
        name: String,
        kind: TypeKind,  // Struct, Enum, Trait
        field_count: usize,
        derive_count: usize,
    },
    /// An import statement.
    Import {
        path: String,
        is_external: bool,  // true = from dependency, false = crate::
    },
    /// A layer assignment (from HDAG node position).
    Layer {
        depth: u8,
        artifact_type: ArtifactType,  // Model, Query, Service, Api, ...
    },
    /// A relationship (FK, trait impl, module dependency).
    Relation {
        source: String,
        target: String,
        kind: RelationKind,  // ForeignKey, Implements, Uses
    },
}

#[derive(Debug, Clone, Serialize, Deserialize, Hash, Eq, PartialEq)]
pub enum TypeKind { Struct, Enum, Trait }

#[derive(Debug, Clone, Serialize, Deserialize, Hash, Eq, PartialEq)]
pub enum ArtifactType {
    Model, Query, Service, Api, Config, Auth, Migration,
    Frontend, Test, Static,
}

#[derive(Debug, Clone, Serialize, Deserialize, Hash, Eq, PartialEq)]
pub enum RelationKind { ForeignKey, Implements, Uses, Contains }
\end{lstlisting}

\subsection{Extraction}

\begin{lstlisting}[style=rust]
/// Extract all Resonites from a Rust source string.
pub fn extract_resonites(
    code: &str, 
    file_path: &str
) -> Vec<Resonite> {
    let parsed = isls_reader::parse_string(code, Language::Rust);
    let mut resonites = Vec::new();
    
    // Functions
    for func in &parsed.functions {
        resonites.push(Resonite::Fn {
            name: func.name.clone(),
            arity: func.params.len(),
            return_type: func.return_type.clone()
                .unwrap_or_default(),
            is_pub: func.is_pub,
            is_async: func.is_async,
        });
    }
    
    // Types
    for typ in &parsed.types {
        resonites.push(Resonite::Type {
            name: typ.name.clone(),
            kind: match typ.kind {
                "struct" => TypeKind::Struct,
                "enum" => TypeKind::Enum,
                "trait" => TypeKind::Trait,
                _ => TypeKind::Struct,
            },
            field_count: typ.fields.len(),
            derive_count: typ.derives.len(),
        });
    }
    
    // Imports
    for import in &parsed.imports {
        resonites.push(Resonite::Import {
            path: import.path.clone(),
            is_external: !import.path.starts_with("crate::"),
        });
    }
    
    // Layer (inferred from file path)
    if let Some(artifact_type) = infer_artifact_type(file_path) {
        resonites.push(Resonite::Layer {
            depth: infer_layer_depth(file_path),
            artifact_type,
        });
    }
    
    resonites
}
\end{lstlisting}

\subsection{Resonite Vectors}

A file's Resonite set is its \textbf{infogenetische Signatur} --- 
the atomic description of what the file \emph{is}, independent of 
its textual content. Two files with identical Resonite sets are 
structurally identical even if their code differs.

The Codematrix (Section~\ref{sec:codematrix}) is computed FROM 
the Resonite set. Resonites are the atoms. The Codematrix is 
the measurement.


% ══════════════════════════════════════════════════════════════════════════════
\section{Multiskalen-Gate}
\label{sec:gates}
% ══════════════════════════════════════════════════════════════════════════════

From the Infogenetik Spektrale Zielsignaturen \S4: three scales 
of gating inserted into the StagedClosure pipeline.

\subsection{Gate Architecture}

\begin{lstlisting}[style=bash]
S4 Solve (topological traversal):
  For each node:
    generate content (structural or LLM/Swarm)
    |
    +-- MIKRO-GATE: check this file in isolation
    |     Pass -> write file, continue
    |     Fail -> Swarm-Coagula or skip with warning
    |
  After each completed layer (all nodes of same depth):
    |
    +-- MESO-GATE: check cross-file consistency in layer
    |     Pass -> continue to next layer
    |     Fail -> identify inconsistent files, re-generate
    |
S5 Gate (existing):
  |
  +-- MAKRO-GATE: cargo check
       Pass -> S7 Emit
       Fail -> S6 Coagula (unchanged)
\end{lstlisting}

\subsection{Mikro-Gate: Per-File Validation}

Applied immediately after LLM/Swarm generates a file, before 
writing to disk.

\begin{lstlisting}[style=rust]
pub struct MikroGateResult {
    pub pass: bool,
    pub codematrix: Codematrix,
    pub violations: Vec<MikroViolation>,
}

pub enum MikroViolation {
    /// File has zero functions (likely garbage output).
    EmptyContent,
    /// Import references a symbol not in ProvidedSymbols.
    UnknownImport(String),
    /// Function name violates naming convention (Rule 3).
    NamingViolation { name: String, expected_pattern: String },
    /// File has no `use` statements (isolated, probably wrong).
    NoImports,
    /// Topology mismatch: file type doesn't match HDAG layer.
    TopologyMismatch { expected: ArtifactType, actual: ArtifactType },
    /// Resonance below threshold.
    LowResonance { value: f64, threshold: f64 },
}

pub fn mikro_gate(
    code: &str,
    node: &HdagNode,
    provided_symbols: &[ProvidedSymbol],
    threshold: f64,
) -> MikroGateResult {
    let resonites = extract_resonites(code, &node.path);
    let codematrix = compute_codematrix(&resonites, node);
    let mut violations = Vec::new();
    
    // Check 1: Non-empty content
    let fn_count = resonites.iter()
        .filter(|r| matches!(r, Resonite::Fn { .. }))
        .count();
    if fn_count == 0 && node.node_type == NodeType::Llm {
        violations.push(MikroViolation::EmptyContent);
    }
    
    // Check 2: All imports resolve to ProvidedSymbols
    for r in &resonites {
        if let Resonite::Import { path, is_external: false } = r {
            if !provided_symbols.iter().any(|ps| 
                path.contains(&ps.import_path)) {
                violations.push(
                    MikroViolation::UnknownImport(path.clone())
                );
            }
        }
    }
    
    // Check 3: Naming convention (Rule 3)
    // ... (entity-specific naming checks)
    
    // Check 4: Codematrix resonance
    let res = codematrix.resonance();
    if res < threshold {
        violations.push(MikroViolation::LowResonance {
            value: res, threshold,
        });
    }
    
    MikroGateResult {
        pass: violations.is_empty(),
        codematrix,
        violations,
    }
}
\end{lstlisting}

Default Mikro-Gate threshold: $\tau_\text{mikro} = 0.15$ 
(very permissive --- only catches obvious garbage).

\textbf{On failure:} If the Swarm produced the answer, try the 
second-best Mandorla candidate. If single-oracle, log a warning 
and write the file anyway (the Makro-Gate will catch it).

\subsection{Meso-Gate: Cross-File Consistency}

Applied after all nodes of a given HDAG layer are processed.

\begin{lstlisting}[style=rust]
pub struct MesoGateResult {
    pub pass: bool,
    pub layer: u8,
    pub file_count: usize,
    pub consistency_score: f64,
    pub violations: Vec<MesoViolation>,
}

pub enum MesoViolation {
    /// Two files in the same layer define the same function.
    DuplicateFunction { name: String, files: Vec<String> },
    /// A service calls a query function that doesn't exist.
    MissingDependency { caller: String, callee: String },
    /// Naming inconsistency across files in the layer.
    InconsistentNaming { files: Vec<String>, pattern: String },
    /// Codematrix variance too high within the layer.
    HighVariance { axis: String, variance: f64 },
}

pub fn meso_gate(
    layer_files: &[(String, String, Codematrix)], // (path, code, cm)
    provided: &[ProvidedSymbol],
) -> MesoGateResult {
    let mut violations = Vec::new();
    
    // Check 1: No duplicate function definitions
    let mut all_fns: HashMap<String, Vec<String>> = HashMap::new();
    for (path, code, _) in layer_files {
        let resonites = extract_resonites(code, path);
        for r in &resonites {
            if let Resonite::Fn { name, .. } = r {
                all_fns.entry(name.clone())
                    .or_default()
                    .push(path.clone());
            }
        }
    }
    for (name, files) in &all_fns {
        if files.len() > 1 {
            violations.push(MesoViolation::DuplicateFunction {
                name: name.clone(),
                files: files.clone(),
            });
        }
    }
    
    // Check 2: Codematrix variance within layer
    if layer_files.len() >= 2 {
        let cms: Vec<&Codematrix> = layer_files.iter()
            .map(|(_, _, cm)| cm).collect();
        for (axis_name, extractor) in &[
            ("R", |cm: &Codematrix| cm.r),
            ("F", |cm: &Codematrix| cm.f),
            ("S", |cm: &Codematrix| cm.s),
        ] {
            let values: Vec<f64> = cms.iter()
                .map(|cm| extractor(cm)).collect();
            let mean = values.iter().sum::<f64>() / values.len() as f64;
            let var = values.iter()
                .map(|v| (v - mean).powi(2))
                .sum::<f64>() / values.len() as f64;
            if var > 0.15 {
                violations.push(MesoViolation::HighVariance {
                    axis: axis_name.to_string(),
                    variance: var,
                });
            }
        }
    }
    
    let consistency = 1.0 - (violations.len() as f64 * 0.2)
        .min(1.0);
    
    MesoGateResult {
        pass: violations.is_empty(),
        layer: layer_files.first()
            .map(|(p, _, _)| infer_layer_depth(p))
            .unwrap_or(0),
        file_count: layer_files.len(),
        consistency_score: consistency,
        violations,
    }
}
\end{lstlisting}

\textbf{On failure:} Log violations. Do NOT re-generate 
(too expensive). Proceed to Makro-Gate. The Meso violations 
are informational in I1 --- they feed the metrics and fitness 
system but don't block emission. In I2 they become blocking.

\subsection{Makro-Gate: Unchanged}

\texttt{cargo check} in S5. No changes. But now it receives 
files that already passed Mikro-Gate, so the failure rate drops.


% ══════════════════════════════════════════════════════════════════════════════
\section{Fitness-R\"uckkopplung}
\label{sec:fitness}
% ══════════════════════════════════════════════════════════════════════════════

\subsection{Norm Fitness}

Every norm in the catalog gets a fitness value:

\begin{lstlisting}[style=rust]
/// Extended norm with fitness tracking.
pub struct NormFitness {
    pub norm_id: String,
    pub fitness: f64,           // [0, 1], default 0.5
    pub activation_count: u64,
    pub success_count: u64,     // Makro-Gate passed
    pub failure_count: u64,     // Makro-Gate failed
    pub last_activated: String,
    pub avg_codematrix: Option<Codematrix>,  // mean 5D of files
                                             // generated under this norm
}
\end{lstlisting}

\subsection{Update Rule}

After every generation (CLI or Cockpit or Swarm):

\begin{equation}
\phi_{t+1}(n) = \alpha \cdot \phi_t(n) + (1 - \alpha) \cdot r_t
\label{eq:fitness}
\end{equation}

where $r_t = 1$ if Makro-Gate passed, $r_t = 0$ if failed, 
$\alpha = 0.9$ (exponential smoothing, recent results matter 
more).

\subsection{Fitness-Weighted Activation}

When \texttt{match\_description()} activates norms, the 
confidence score is multiplied by fitness:

\begin{equation}
\text{score}_\text{adjusted}(n) = \text{score}_\text{keyword}(n) \cdot \phi(n)
\end{equation}

A norm with $\phi = 0.3$ (mostly fails) gets suppressed even 
if the keyword match is strong. A norm with $\phi = 0.95$ 
(almost always succeeds) gets boosted.

\subsection{Storage}

Fitness data stored in \texttt{\textasciitilde/.isls/fitness.json}, 
alongside \texttt{norms.json}. Updated after every generation. 
Loaded at startup by \texttt{NormRegistry}.

\subsection{CLI Inspection}

\begin{lstlisting}[style=bash]
isls norms fitness

ISLS Norm Fitness Report
---
ISLS-NORM-0042 (CRUD-Entity):     0.95  (19/20 success)
ISLS-NORM-0088 (JWT-Auth):        0.90  (18/20)
ISLS-NORM-0096 (Pagination):      0.95  (19/20)
ISLS-NORM-INFRA-WEB:              0.92  (12/13)
ISLS-NORM-AUTO-0001 (Auto-Pattern): 0.85  (6/7)

Weakest:
ISLS-NORM-0112 (Inventory):       0.50  (1/2) -- insufficient data
\end{lstlisting}


% ══════════════════════════════════════════════════════════════════════════════
\section{Kontraktivit\"ats-Monitor}
\label{sec:contraction}
% ══════════════════════════════════════════════════════════════════════════════

From Infogenetik Spektrale Zielsignaturen \S3: the Solve-Coagula 
stack is provably convergent if the operator is contractive 
(Lipschitz $L < 1$). In practice: each Coagula cycle should 
reduce the error count.

\begin{lstlisting}[style=rust]
pub struct ContractionMonitor {
    pub error_counts: Vec<usize>,  // per Coagula cycle
    pub is_contractive: bool,
}

impl ContractionMonitor {
    pub fn record_cycle(&mut self, errors: usize) {
        self.error_counts.push(errors);
        if self.error_counts.len() >= 2 {
            let prev = self.error_counts[self.error_counts.len() - 2];
            let curr = self.error_counts[self.error_counts.len() - 1];
            self.is_contractive = curr < prev;
        }
    }
    
    /// Contraction ratio: < 1.0 means converging.
    pub fn contraction_ratio(&self) -> Option<f64> {
        if self.error_counts.len() >= 2 {
            let prev = self.error_counts[self.error_counts.len() - 2] as f64;
            let curr = self.error_counts[self.error_counts.len() - 1] as f64;
            if prev > 0.0 {
                Some(curr / prev)
            } else {
                Some(0.0) // No errors = converged
            }
        } else {
            None
        }
    }
}
\end{lstlisting}

Logged in \texttt{GenerationMetrics}:

\begin{lstlisting}[style=rust]
pub struct GenerationMetrics {
    // ... existing fields ...
    pub contraction_ratios: Vec<f64>,
    pub was_contractive: bool,
    pub codematrix_avg: Option<Codematrix>,
    pub mikro_gate_pass_rate: f64,
    pub meso_gate_pass_rate: f64,
}
\end{lstlisting}

Displayed in Studio (Erschaffen mode) after forge:

\begin{lstlisting}[style=bash]
  | Kontraktivitaet: 
  |   Cycle 1: 12 errors (initial)
  |   Cycle 2: 3 errors  (ratio: 0.25 -- contractive)
  |   Cycle 3: 0 errors  (converged)
  |
  | 5D Codematrix (project average):
  |   R=0.62  F=0.78  T=0.91  S=0.88  E=0.55
  |   Resonance: 0.74
\end{lstlisting}


% ══════════════════════════════════════════════════════════════════════════════
\section{Work Packages}
% ══════════════════════════════════════════════════════════════════════════════

\subsection{W1: Resonite + Codematrix Module}

\subsubsection{New Files}

\begin{longtable}{lp{10cm}}
\toprule
\textbf{File} & \textbf{Purpose} \\
\midrule
\texttt{isls-forge-llm/src/codematrix.rs}
  & \texttt{Codematrix} struct, 5 axis computation functions, 
    \texttt{Resonite} enum, \texttt{extract\_resonites()}, 
    \texttt{compute\_codematrix()}. \\
\bottomrule
\end{longtable}

\subsubsection{Dependencies}

Uses \texttt{isls\_reader::parse\_string()} for code parsing. 
No new external dependencies.

\subsubsection{W1 Verification}

\begin{enumerate}[nosep]
\item Unit test: extract resonites from a known Rust file, 
      verify function/struct/import counts.
\item Unit test: compute codematrix for a query file, verify 
      axes are in $[0, 1]$.
\item Unit test: codematrix distance between identical files $= 0$.
\item Unit test: codematrix resonance for a well-formed file $> 0.5$.
\end{enumerate}


\subsection{W2: Mikro-Gate + Meso-Gate}

\subsubsection{New Files}

\begin{longtable}{lp{10cm}}
\toprule
\textbf{File} & \textbf{Purpose} \\
\midrule
\texttt{isls-forge-llm/src/gates.rs}
  & \texttt{mikro\_gate()}, \texttt{meso\_gate()}, violation 
    types, gate result structs. \\
\bottomrule
\end{longtable}

\subsubsection{Modified Files}

\begin{longtable}{lp{10cm}}
\toprule
\textbf{File} & \textbf{Change} \\
\midrule
\texttt{isls-forge-llm/src/staged\_closure.rs}
  & After each LLM node: call \texttt{mikro\_gate()}. 
    After each completed layer: call \texttt{meso\_gate()}. 
    Log results. Store codematrix per file. \\
\texttt{isls-forge-llm/src/lib.rs}
  & \texttt{pub mod codematrix; pub mod gates;} \\
\bottomrule
\end{longtable}

\subsubsection{W2 Verification}

\begin{enumerate}[nosep]
\item Run \texttt{isls forge-chat -m "Pet shop"} --- verify 
      Mikro-Gate log lines appear for each LLM node.
\item Verify Meso-Gate log lines appear after each layer.
\item All 6 D4 pipeline apps still compile (backward compat).
\end{enumerate}


\subsection{W3: Fitness-R\"uckkopplung}

\subsubsection{New Files}

\begin{longtable}{lp{10cm}}
\toprule
\textbf{File} & \textbf{Purpose} \\
\midrule
\texttt{isls-norms/src/fitness.rs}
  & \texttt{NormFitness}, \texttt{FitnessStore}, 
    \texttt{update\_fitness()}, \texttt{load/save\_fitness()},
    fitness-weighted activation. \\
\bottomrule
\end{longtable}

\subsubsection{Modified Files}

\begin{longtable}{lp{10cm}}
\toprule
\textbf{File} & \textbf{Change} \\
\midrule
\texttt{isls-norms/src/lib.rs}
  & \texttt{pub mod fitness;} Integrate fitness into 
    \texttt{match\_description()} scoring. \\
\texttt{isls-forge-llm/src/staged\_closure.rs}
  & After S7 Emit or S6 exhaustion: call 
    \texttt{update\_fitness(norm\_ids, success)}. \\
\texttt{isls-cli/src/cmd\_norms.rs}
  & Add \texttt{isls norms fitness} subcommand. \\
\bottomrule
\end{longtable}

\subsubsection{W3 Verification}

\begin{enumerate}[nosep]
\item Run 3 successful generations. Check 
      \texttt{\textasciitilde/.isls/fitness.json} has entries 
      with $\phi > 0.5$.
\item Manually set a norm's fitness to 0.1. Verify it is 
      suppressed in \texttt{match\_description()} output.
\item \texttt{isls norms fitness} shows the fitness table.
\end{enumerate}


\subsection{W4: Kontraktivit\"ats-Monitor + Metrics}

\subsubsection{Modified Files}

\begin{longtable}{lp{10cm}}
\toprule
\textbf{File} & \textbf{Change} \\
\midrule
\texttt{isls-forge-llm/src/staged\_closure.rs}
  & Track error count per Coagula cycle. Log contraction 
    ratio. \\
\texttt{isls-forge-llm/src/metrics.rs}
  & Add \texttt{contraction\_ratios}, 
    \texttt{was\_contractive}, \texttt{codematrix\_avg}, 
    \texttt{mikro\_gate\_pass\_rate}, 
    \texttt{meso\_gate\_pass\_rate} to 
    \texttt{GenerationMetrics}. \\
\bottomrule
\end{longtable}

\subsubsection{W4 Verification}

\begin{enumerate}[nosep]
\item Trigger a Coagula-producing generation. Verify 
      contraction ratios logged in metrics.jsonl.
\item \texttt{isls metrics --last 1} shows 5D codematrix 
      and contraction data.
\end{enumerate}


% ══════════════════════════════════════════════════════════════════════════════
\section{Constraints for Claude Code}
% ══════════════════════════════════════════════════════════════════════════════

\begin{enumerate}
\item \textbf{Resonite extraction} uses \texttt{isls\_reader::parse\_string()}. 
      Do NOT reimplement Rust parsing.
\item \textbf{Codematrix values} MUST be in $[0, 1]$. Clamp all results.
\item \textbf{Mikro-Gate failures are warnings}, not errors. They 
      do NOT prevent file writing in I1. They feed metrics.
\item \textbf{Meso-Gate failures are informational} in I1. 
      They log but don't block.
\item \textbf{Fitness persistence} uses \texttt{\textasciitilde/.isls/fitness.json}. 
      Same directory as \texttt{norms.json}. Same 
      \texttt{dirs\_path()} with USERPROFILE fallback.
\item \textbf{Fitness default} for new norms: $\phi = 0.5$ 
      (neutral, neither boosted nor suppressed).
\item \textbf{Do not modify} isls-merkaba, isls-gateway, isls-chat, 
      isls-reader, isls-code-topo, isls-hypercube, isls-types.
\item \textbf{All existing tests MUST pass.}
\item Branch: \texttt{claude/implement-isls-i1-spec}.
\item Commit format: \texttt{I1/W\{n\}: <description>}.
\end{enumerate}


% ══════════════════════════════════════════════════════════════════════════════
\section{Dependency Graph}
% ══════════════════════════════════════════════════════════════════════════════

\begin{verbatim}
W1 (Resonite + Codematrix)
  |
  +---> W2 (Mikro-Gate + Meso-Gate)
  |
  +---> W3 (Fitness-Rueckkopplung)
          |
          v
        W4 (Kontraktivitaets-Monitor + Metrics)
\end{verbatim}

W1 first (foundation). W2 and W3 depend on W1 independently. 
W4 integrates both into metrics.

Recommended: W1 $\to$ W2 $\to$ W3 $\to$ W4.


% ══════════════════════════════════════════════════════════════════════════════
\section{Acceptance Criteria}
% ══════════════════════════════════════════════════════════════════════════════

\begin{longtable}{clp{9cm}}
\toprule
\textbf{\#} & \textbf{Pass} & \textbf{Criterion} \\
\midrule
1 & & \texttt{cargo build --workspace} compiles. \\
2 & & \texttt{cargo test --workspace} passes. \\
3 & & Codematrix unit tests pass (4+ tests). \\
4 & & Mikro-Gate log lines visible in forge output. \\
5 & & Meso-Gate log lines visible after layer completion. \\
6 & & \texttt{fitness.json} created after first generation. \\
7 & & \texttt{isls norms fitness} shows fitness table. \\
8 & & Contraction ratios in \texttt{metrics.jsonl}. \\
9 & & Pipeline bat: 6/6 PASS (backward compat). \\
\bottomrule
\end{longtable}

I1 is passed when criteria 1--3, 6, and 9 are fulfilled. 
Criteria 4--5 and 7--8 require a generation run with LLM.


\vfill
\begin{center}
\color{isls-gray}\rule{0.5\textwidth}{0.4pt}\\[0.5em]
{\small End of I1 Specification.\\[0.3em]
Die Norm ist nicht der Code.\\
Die Norm ist die Erkenntnis \"uber Code.\\
Und jetzt messen wir die Erkenntnis in f\"unf Dimensionen.}
\end{center}

\end{document}
